\documentclass[11pt, a4paper, british]{article}

\usepackage[style=apa, backend=biber]{biblatex}
\usepackage[margin=1in]{geometry}  
\usepackage{setspace}
\usepackage{wrapfig}
\usepackage{mathrsfs}
\usepackage{boxedminipage}
\usepackage{pifont}
\usepackage{euscript}
\usepackage{boxedminipage}
\usepackage{dcolumn}
\usepackage{float}
\usepackage{fancybox}
\usepackage{fancyhdr}
\usepackage{rotating}
\usepackage{multirow}
\usepackage{latexsym}
\usepackage{amssymb}
\usepackage{amsmath}
\usepackage{epsfig}
\usepackage{color}
\usepackage{verbatim}
\usepackage{subfigure}
\usepackage{mathrsfs}
\usepackage{url}
\usepackage{lscape}
\usepackage{psfrag}
\usepackage{graphicx}
\usepackage{mparhack} % Fix \margref
\usepackage[ruled,vlined,linesnumbered]{algorithm2e}
\usepackage{dsfont}
\usepackage{hyperref}
\usepackage{amsthm}
\usepackage{libertine}
\usepackage[utf8]{inputenc}
\usepackage[margin=1in]{geometry}
\usepackage{multicol}
\usepackage[shortlabels]{enumitem}
\usepackage{siunitx}
\usepackage{cancel}
\usepackage{pgfplots}
\usepackage{tikz}
\usepackage{bbm}
\usepackage{tcolorbox}
\usepackage{xcolor}
\usepackage{booktabs}
\usepackage[utf8]{inputenc}
\usepackage[T1]{fontenc}
\usepackage{libertine}
\usepackage{microtype}
\usepackage{lmodern}
\usepackage{breqn}
\usepackage{nicefrac}
\usepackage{varwidth}
\usepackage{amsmath} % Required for advanced math features
\usepackage{amssymb} % Additional math symbols
\usepackage{mathrsfs} % For \mathcal and \mathscr fonts
\usepackage{tabularx}
\usepackage{booktabs}      % For \toprule, \midrule, \bottomrule
\usepackage{tabularx}      % For auto-sized columns (X)
\usepackage{colortbl}      % Extra safety: ensures \rowcolor, etc., work
\usepackage{array}         % Often helps with column types
\usepackage{caption} 
\usepackage{tikz}
\usepackage{calc}
\usepackage{threeparttable}
\usepackage{titling}
\usetikzlibrary{arrows.meta, positioning, decorations.pathreplacing, shapes.symbols}

\pgfplotsset{compat=1.18}
\definecolor{ltgray}{HTML}{F2F2F2}
\definecolor{mdgray}{HTML}{E3E3E3}
\definecolor{mygreen}{rgb}{0,0.6,0}
\definecolor{mygray}{rgb}{0.5,0.5,0.5}
\definecolor{mymauve}{rgb}{0.58,0,0.82}
\definecolor{applegreen}{rgb}{0.55, 0.41, 0.0}
\definecolor{gray1}{rgb}{0.98, 0.98, 0.98}
\definecolor{black}{rgb}{0, 0, 0}
\definecolor{antique1}{rgb}{0.98, 0.92, 0.84}
\definecolor{cream}{rgb}{1.0, 0.99, 0.82}
\definecolor{darkpink}{rgb}{0.91, 0.33, 0.5}
\definecolor{darkorchid}{rgb}{0.6, 0.2, 0.8}
\definecolor{indigo}{rgb}{0.29, 0.0, 0.51}
\definecolor{islamicgreen}{rgb}{0.0, 0.56, 0.0}
\definecolor{vlgray}{gray}{0.9}
\definecolor{lgray}{gray}{0.7}
\definecolor{defncolor}{rgb}{0.76,0.96,0.95}
\definecolor{examplecolor}{rgb}{0.76, 0.96, 0.76} % Light green background color
\definecolor{headercolor}{rgb}{0.2,0.2,0.2} % Dark header color
\definecolor{appcolor}{rgb}{0.95, 0.80, 0.78} 
\definecolor{darkblue}{rgb}{0.0, 0.0, 0.5} % Dark blue color
\definecolor{darkgreen}{rgb}{0.0, 0.5, 0.0} % Dark green color
\definecolor{darkred}{rgb}{0.5, 0.0, 0.0} % Dark red color



\newcommand{\myurl}[1]{{\textcolor{blue}{\url{#1}}}}
\newcommand{\myhref}[2]{\textcolor{blue}{{\href{#1}{#2}}}}
\newcommand{\preface}[1][Preface]{\chapter*{#1}\markboth{#1}{#1}}
\newcommand{\code}[1]{{\normalfont\ttfamily\hyphenchar\font=-1 #1}}
\newcommand{\pkg}[1]{{\fontseries{b}\selectfont #1}}
\newcommand{\str}[1]{{{\tt \textquotesingle}}{\ttfamily#1}{\tt \textquotesingle}}

%commands
\newcommand{\bb}[1]{\mathbb{#1}}
\newcommand{\X}{\textcolor{blue}{\mathbf{X}}} %model matrix
\newcommand{\Y}{\textcolor{blue}{\mathbf{Y}}} %output matrix
\newcommand{\indsim}{\stackrel{\mathrm{ind}}{\sim}}
\newcommand{\simind}{\indsim}
\newcommand{\simindt}{{\:{\sim}_{\mathrm{ind}}\:}}
\newcommand{\tilda}{\mathem{\sim}}
\newcommand{\blar}{\textcolor{black}>}
\newcommand{\plar}{\textcolor{black}{>>}}
\newcommand{\blpl}{\textcolor{black}+}
\newcommand{\sgn}{\mathrm{sgn}}
\newcommand{\la}{\leftarrow}
\newcommand{\st}{\: : \:} % such that
\newcommand {\otoprule }{\midrule[\heavyrulewidth]}
\newcommand{\Span}[1]{\mathrm{Span}\left\{ #1 \right\}}
\newcommand{\ceil}[1]{\left\lceil #1 \right\rceil}
\newcommand{\floor}[1]{\left\lfloor #1 \right\rfloor}
\newcommand{\dd}{\partial}
\newcommand{\cZ}{\mathcal{Z}}
\newcommand{\Shan}{{\EuScript{H}}}
\newcommand{\KL}{{\EuScript{D}}}
\newcommand{\El}{{\EuScript{E}}}
\newcommand{\MI}{\EuScript{M}} % mutual information
\newcommand{\Lh}{\EuScript{L}} % likelihood
\newcommand{\Fish}{\EuScript{I}}
\newcommand{\Sc}{\EuScript{S}} % score function
\newcommand{\bigO}{\c O}
\newcommand{\btau}{\vect{\tau}}
\newcommand{\bphi}{\vect{\phi}}
\newcommand{\bbeta}{\vect{\beta}}
\newcommand{\bepsilon}{\vect{\epsilon}}
\newcommand{\bomega}{\vect{\omega}}
\newcommand{\bmu}{\vect{\mu}}
\newcommand{\bsig}{\vect{\sigma}}
\newcommand{\lito}{$o$}
\renewcommand{\v}[1]{\boldsymbol{#1}}
\newcommand{\bS}{\v{S}}
\newcommand{\bL}{\v{L}}
\newcommand{\bl}{\mbox{\boldmath $\ell$}}
\newcommand{\bs}{\v{s}}
\newcommand{\bQ}{\v{Q}}
\newcommand{\bG}{\v{G}}
\newcommand{\bg}{\v{g}}
\newcommand{\bSigma}{\mbox{\boldmath $\Sigma$}}
\newcommand{\bB}{\v{B}}
\newcommand{\bz}{\v{z}}
\newcommand{\bincof}[2]{{{#1} \choose {#2}}}


\newcommand*\xbar[1]{%
   \hbox{%
     \vbox{%
       \hrule height 0.8pt % The actual bar
       \kern0.3ex%         % Distance between bar and symbol
       \hbox{%
         \kern+0em%      % Shortening on the left side
         \ensuremath{#1}%
         \kern-0.0em%      % Shortening on the right side
       }%
     }%
   }%
}

\newcommand{\T}{{\top}}
\newcommand{\gr}{\mbox{$\nabla$}}
\newcommand{\bK}{\v{K}}
\newcommand{\bX}{\v{X}}
\newcommand{\bx}{\v{x}}
\newcommand{\bk}{\v{k}}
\newcommand{\bq}{\v{q}}
\newcommand{\be}{\v{e}}
\newcommand{\bF}{\v{F}}
\newcommand{\br}{\v{r}}
\newcommand{\by}{\v{y}}
\newcommand{\bZ}{\v{Z}}
\newcommand{\bu}{\v{u}}
\newcommand{\bU}{\v{U}}
\newcommand{\bY}{\v{Y}}
\newcommand{\bP}{\v{P}}
\newcommand{\bR}{\v{R}}
\newcommand{\bV}{\v{V}}
\newcommand{\bW}{\v{W}}
\newcommand{\bT}{\v{T}}
\newcommand{\bt}{\v{t}}

\newcommand{\dotsim}{\stackrel{\centerdot}{\sim}}
\renewcommand{\le}{\leq}
\renewcommand{\ge}{\geq}
\renewcommand{\tilde}{\widetilde}
\newcommand{\bnu}{\vect{\nu}}
\newcommand{\btheta}{\vect{\theta}}
\newcommand{\blambda}{\vect{\lambda}}
\newcommand{\balpha}{\vect{\alpha}}
\newcommand{\halmos}{\vspace{3mm} \hfill \mbox{$\Box$}}

\newcommand{\bE}{\v{E}}
\newcommand{\ba}{\v{a}}
\newcommand{\bh}{\v{h}}
\newcommand{\bfb}{\v{b}}
\newcommand{\bA}{\v{A}}
\newcommand{\bM}{\v{M}}
\newcommand{\bN}{\v{N}}
\newcommand{\bc}{\v{c}}
\newcommand{\bbf}{\v{f}}
\newcommand{\bd}{\v{d}}
\newcommand{\bD}{\v{D}}
\newcommand{\bw}{\v{w}}

\newcommand{\Var}{\mathbb{V}\mathrm{ar}}
\newcommand{\E}{\mathbb{E}}
\newcommand{\var}{\Var}

\newcommand{\Cov}{\mathbb{C}\mathrm{ov}}
\newcommand{\cov}{\Cov}

\newcommand{\Supp}{\mathrm{Supp}}
\newcommand{\matlab}{\mathrm{M}\mathrm{{\scriptstyle ATLAB}}}


\newcommand{\bp}{\v{p}}
\newcommand{\bC}{\v{C}}
\newcommand{\bii}{\v{i}}
\newcommand{\bjj}{\v{j}}
\newcommand{\bII}{\v{I}}
\newcommand{\bv}{\v{v}}

\renewcommand{\epsilon}{\varepsilon}
% \renewcommand{\rho}{\varrho}
\renewcommand{\log}{\ln}
\renewcommand{\hat}{\widehat}
\renewcommand{\leq}{\leqslant}
\renewcommand{\geq}{\geqslant}

\newcommand{\argmax}{\mathop{\rm argmax}}
\newcommand{\argmin}{\mathop{\rm argmin}}

\newcommand{\iid}{\text{iid }}

%distributions
% Bernoulli distribution
\newcommand{\Ber}{{\sf Ber}}
\newcommand{\ber}{\Ber}

% Erlang distribution
\newcommand{\Erl}{{\sf Erl}}
\newcommand{\erl}{\Erl}

% Binomial distribution
\newcommand{\Bin}{{\sf Bin}}
\newcommand{\bin}{\Bin}

% Cauchy distribution
\newcommand{\Cauchy}{{\sf Cauchy}}
\newcommand{\cauchy}{\Cauchy}

% Negative binomial distribution
\newcommand{\NegBin}{{\sf NegBin}}
\newcommand{\negbin}{\NegBin}
\newcommand{\negBin}{\NegBin}
\newcommand{\Negbin}{\NegBin}
\newcommand{\nbin}{\NegBin}
\newcommand{\Nbin}{\NegBin}
\newcommand{\NBin}{\NegBin}

% Multinomial distribution
\newcommand{\Mnom}{{\sf Mnom}}
\newcommand{\mnom}{\Mnom}

% Geometric distribution
\newcommand{\Geo}{{\sf Geom}}
\newcommand{\geo}{\Geo}
\newcommand{\Geom}{\Geo}
\newcommand{\geom}{\Geo}
\newcommand{\G}{\Geo}
\newcommand{\NE}{{\sf NE}}

% Hypergeometric distribution
\newcommand{\Hyp}{{\sf Hyp}}

% Poisson distribution
\newcommand{\Poi}{{\sf Poi}}
\newcommand{\poi}{\Poi}
\newcommand{\Po}{\Poi}
\newcommand{\po}{\Poi}

% Uniform distribution (continuous)
%\newcommand{\U}{{\sf Unif}}
\newcommand{\U}{\EuScript{U}}
% Exponential distribution
\newcommand{\Ex}{{\sf Exp}}
\newcommand{\ex}{\Ex}

% Normal / Gaussian distribution
\newcommand{\Nor}{\EuScript{N}}
\newcommand{\nor}{\Nor}

% Pareto distribution
\newcommand{\Pareto}{{\sf Pareto}}
\newcommand{\pareto}{\Pareto}
\newcommand{\ParetoI}{{\sf ParetoI}}

% Negative hypergeometric distribution
\newcommand{\NegHyp}{{\sf NegHyp}}

% Phase-type distributions
\newcommand{\DiscPhase}{{\sf DPH}}
\newcommand{\ContPhase}{{\sf PH}}

% F distribution
\newcommand{\Fdist}{{\sf F}}

% Extreme value distributions
\newcommand{\EVI}{{\sf ExValueI}}
\newcommand{\EVII}{{\sf ExValueII}}
\newcommand{\EVIII}{{\sf ExValueIII}}
\newcommand{\Gumbel}{{\sf Gumbel}}
\newcommand{\Frechet}{\text{\sf Fr\'{e}chet}}
\newcommand{\gumbel}{\Gumbel}
\newcommand{\frechet}{\Frechet}

% Inverse Gaussian distribution
\newcommand{\IGauss}{{\Wald}}
\newcommand{\InvGauss}{\Wald}
\newcommand{\Wald}{{\sf Wald}}

% Logistic distribution
\newcommand{\Logistic}{{\sf Logistic}}
\newcommand{\logistic}{\Logistic}

% Lognormal distribution
\newcommand{\LogN}{{\sf LogN}}

% Rayleigh distribution
\newcommand{\Ray}{{\sf Rayleigh}}

% Levy \alpha stable distribution
\newcommand{\Stable}{{\sf Stable}}

% Student's t distribution
\newcommand{\Student}{{\sf t}}
\newcommand{\student}{\Student}

% von Mises distribution
\newcommand{\vonMises}{{\sf vonMises}}

% Dirichlet distribution
\newcommand{\Dirichlet}{{\sf Dirichlet}}
\newcommand{\dirichlet}{\Dirichlet}
\newcommand{\Dir}{\Dirichlet}
\newcommand{\dir}{\Dirichlet}

% Gamma distribution
\newcommand{\Gam}{{\sf Gamma}}
\newcommand{\gam}{\Gam}

% Inverse Gamma
\newcommand{\invgam}{{\sf InvGamma}}
% Beta distribution
\newcommand{\Bet}{{\sf Beta}}
\newcommand{\bet}{\Bet}

% Weibull distribution
\newcommand{\Weib}{{\sf Weib}}
\newcommand{\weib}{\Weib}

% Discrete uniform distribution
\newcommand{\DU}{{\sf DU}}

% Shifted exponential distribution
\newcommand{\SE}{{\sf SExp}}

% Double exponential distribution
%\newcommand{\DE}{{\sf DExp}}
\newcommand{\DE}{{\sf Laplace}}
\newcommand{\DExp}{\DE}
\newcommand{\Laplace}{\DE}
\newcommand{\laplace}{\DE}

% Arcsine distribution
\newcommand{\Arcsine}{{\sf Arcsine}}

% Levy distribution
\newcommand{\Levy}{\text{\sf L\'{e}vy}}

% Generic distribution
\newcommand{\Dist}{{\sf Dist}}

% Wishart distribution
\newcommand{\Wishart}{{\sf Wishart}}


\newcommand{\Em}{{\mathbb E}}
\newcommand{\Pm}{{\mathbb P}}  
\newcommand{\R}{{\mathbb R}}
\newcommand{\Rbar}{\overline{\mathbb R}}

\newcommand{\B}{{\cal B}}
\newcommand{\cB}{{\cal B}}
\newcommand{\cE}{{\cal E}}

\newcommand{\F}{\mathbb F}
\newcommand{\N}{\mathbb N}
\newcommand{\Z}{\mathbb Z}
\newcommand{\Q}{\mathbb Q}
\newcommand{\C}{\mathbb C}

\newcommand{\scA}{\mathscr{A}}
\newcommand{\scB}{\mathscr{B}}
\newcommand{\scC}{\mathscr{C}}
\newcommand{\scD}{\mathscr{D}}
\newcommand{\scF}{\mathscr{F}}
\newcommand{\scG}{\mathscr{G}}
\newcommand{\scH}{\mathscr{H}}
\newcommand{\scI}{\mathscr{I}}
\newcommand{\scJ}{\mathscr{J}}
\newcommand{\scL}{\mathscr{L}}
\newcommand{\scM}{\mathscr{M}}
\newcommand{\scP}{\mathscr{P}}
\newcommand{\scS}{\mathscr{S}}
\newcommand{\scR}{\mathscr{R}}
\newcommand{\scU}{\mathscr{U}}
\newcommand{\scE}{\mathscr{E}}
\newcommand{\scT}{\mathscr{T}}
\newcommand{\scV}{\mathscr{V}}
\newcommand{\scW}{\mathscr{W}}
\newcommand{\scX}{\mathscr{X}}
\newcommand{\scY}{\mathscr{Y}}
\newcommand{\scZ}{\mathscr{Z}}

\newcommand{\gvn}{\,|\,}

\newcommand{\e}{\mathrm{e}}

\newcommand{\ds}{\displaystyle}

\newcommand{\vect}[1]{\boldsymbol #1}

\newcommand{\convdistr}{\stackrel{d}{\longrightarrow}}
\newcommand{\distrconv}{\convdistr}
\newcommand{\distreq}{\stackrel{d}{=}}
\newcommand{\eqdistr}{\distreq}

\newcommand{\di}{\mathrm{d}}
\newcommand{\tr}{\mathrm{tr}}

\newcommand{\utterance}[1]{\begin{list}{}%
{\leftmargin\parindent\rightmargin0pt\listparindent0pt\parsep1ex%
\labelwidth0pt\labelsep0pt}\item[]\small\sf$\maltese$~#1 \end{list}}

\newcommand{\cF}{{{\cal F}}}
\newcommand{\cG}{{{\cal G}}}
\newcommand{\cH}{{\cal H}}

\newcommand{\Prob}{\Pm}

\newcommand{\cas}{{\stackrel{\mathrm{a.s.}}{\longrightarrow}\: }}

\newcommand{\ra}{\rightarrow}
\newcommand{\lra}{\leftrightarrow}
\newcommand{\rar}{\rightarrow}
\newcommand{\Rar}{\Rightarrow}
\newcommand{\dar}{\downarrow}

\newcommand{\chg}{\marginpar{*}}

\newcommand{\Ito}{It\^{o}}

\def\eps{\varepsilon}

\newcommand{\matX}{{\cal X}}

\newcommand{\I}{\mathds{1}}
\newcommand{\ii}{\mathrm{i}}

\newcommand{\appr}{\mbox{\footnotesize appr.}}
\newcommand{\approxsim}{\stackrel{\mathrm{approx.}}{\sim}}

\newcommand{\iidsim}{\stackrel{\mathrm{iid}}{\sim}}
\newcommand{\simiid}{\iidsim}
\newcommand{\simiidt}{{\:{\sim}_{\mathrm{iid}}\:}}
\newcommand{\diag}{\mathrm{diag}}

\newcommand{\sech}{\mathrm{sech}}

\newcommand{\chk}[1]{}
\newcommand{\ver}{}

\newcommand{\idef}{\stackrel{\mathrm{def}}{=}}
%
%% Correlation
\newcommand{\Corr}{\mathrm{Corr}}
\newcommand{\corr}{\Corr}

\newcommand{\bm}{\v{m}}

%% Space C^k
\newcommand{\eC}{{\EuScript{C}}}
%

\RequirePackage{xcolor}
\definecolor{beaubleu}{rgb}{0.26,0.31,0.61}
\definecolor{beauvert}{rgb}{0.27,0.52,0.42}
\definecolor{datacol}{rgb}{0.43, 0.21, 0.1}
\definecolor{datacol}{rgb}{0.8, 0.33, 0.0}
%\definecolor{packcol}{rgb}{0.2, 0.33, 0.13}
\definecolor{packcol}{rgb}{0.09, 0.45, 0.27}
\definecolor{methcol}{rgb}{0.16, 0.32, 0.75}
%DK: See http://latexcolor.com/
%
\usepackage{cprotect}
\RequirePackage[formats]{listings}\RequirePackage{textcomp}
\RequirePackage[formats]{listings}\RequirePackage{textcomp}
\newcommand{\file}[1]{\texttt{#1}}
%\newcommand{\dataset}[1]{\textcolor{datacol}{\bf \texttt{#1}}}
\newcommand{\dataset}[1]{\texttt{#1}}
\newcommand{\datafr}[1]{\textcolor{datacol}{\bf \texttt{#1}}}
\newcommand{\packagecap}[1]{\textcolor{packcol}{\bf \texttt{#1}}}% only in captions
\newcommand{\package}[1]{\lstset{basicstyle=\bfseries\ttfamily\color{packcol}}\lstinline{#1}}
\newcommand{\method}[1]{\lstset{basicstyle=\bfseries\ttfamily\color{methcol}}\lstinline{#1}}
\newcommand{\atrr}[1]{\lstset{basicstyle=\bfseries\ttfamily\color{methcol}}\lstinline{#1}}
\newcommand{\key}[1]{{\normalfont\texttt{\uppercase{#1}}}}
\newcommand{\menu}[1]{\texttt{#1}}
\newcommand{\menuentry}[1]{\texttt{#1}}

\definecolor{outcol}{gray}{0.9}
%\definecolor{incol}{rgb}{0.82, 0.94, 0.75}
\definecolor{incol}{rgb}{0.80, 0.91, 0.98}

\lstnewenvironment{Pin}{%
  \lstset{backgroundcolor=\color{incol},
  belowskip=0pt,
  mathescape,
  xleftmargin = 3pt,
  xrightmargin = 3pt,
  frame=single,
  framerule=0pt,
  basicstyle=\footnotesize\ttfamily,
  columns=fixed}}{}
  
\lstnewenvironment{Pout}{%
  \lstset{backgroundcolor=\color{outcol},
  aboveskip= -2pt,  
  xleftmargin = 4pt,
  xrightmargin = 4pt,
  frame=single,
  upquote,
  framerule=1pt,
  basicstyle=\footnotesize\ttfamily,
  columns=fixed}}{}


%\newenvironment{Pin}
%{\VerbatimEnvironment
%\begin{Verbatim}[formatcom=\color{beaubleu},commandchars=\\\{\},fontfamily=courier,fontsize=\small,fontseries=b]}
%{\end{Verbatim}}

%% \newenvironment{Pout}
%% {\vspace*{-0.5cm}
%% \VerbatimEnvironment
%% \begin{Verbatim}[formatcom=\color{beauvert},commandchars=\\\{\},fontfamily=courier,fontsize=\small,fontseries=b,fontshape=it]}
%% {\end{Verbatim}}



\newenvironment{Pcode}[1]{%
  \tcblisting{breakable,listing only,colback=black!1!cream, colframe=red!5!black, colbacktitle=black, enhanced,
 		 attach boxed title to top
    center={xshift=-3cm,yshift=-2mm},arc = 0mm,title={ {\ttfamily\large #1}},
            	   listing options={frame=line,language=Python,
basicstyle=\footnotesize\ttfamily,        % the size of the fonts that are used for the code
  numbers=left,                   % where to put the line-numbers
  numberstyle=\tiny\color{blue},  % the style that is used for the
                                % line-numbers
  xleftmargin=5pt,
  aboveskip=0pt,
columns=fixed,
basewidth=6pt,
  stepnumber=1,                   % the step between two line-numbers. If it's 1, each line
                                  % will be numbered
  numbersep=7pt,                  % how far the line-numbers are from
                                % the code
backgroundcolor=\color{cream},  % choose the background color. You must add \usepackage{color}
  showspaces=false,               % show spaces adding particular underscores
  showstringspaces=false,         % underline spaces within strings
  showtabs=false,                 % show tabs within strings adding particular underscores
  %frame=single,                   % adds a frame around the code
  %frame=shadowbox,
  rulecolor=\color{black},        % if not set, the frame-color may be changed on line-breaks within not-black text (e.g. commens (green here))
  tabsize=2,                     % sets default tabsize to 2 spaces
  captionpos=b,                  % sets the caption-position to bottom
  breaklines=true,                % sets automatic line breaking
  breakatwhitespace=false,        % sets if automatic breaks should only happen at whitespace
  title=\lstname,                 % show the filename of files included with \lstinputlisting;
                                  % also try caption instead of title
  keywordstyle=\color{blue},      % keyword style
  commentstyle=\color{islamicgreen},   % comment style
  stringstyle=\color{red},      % string literal style
 % escapeinside={\%*}{*)},         % if you want to add a comment within your code
  morekeywords={*,...}
  			 }
    		}
    	}
{\endtcblisting}


\newenvironment{PC}{%
\tcblisting{breakable,listing only,colback=cream, enhanced, sharpish
corners,  boxrule=1pt,after, listing options =
   {language=Python,
  basicstyle=\footnotesize\ttfamily,
  numbers=none,
  backgroundcolor=\color{cream},
  xleftmargin=-10pt,
  aboveskip=-4pt,
  belowskip=-6pt,
  columns=fixed,
  basewidth=6pt,
  showspaces=false,
  showstringspaces=false,
  showtabs=false,
  tabsize=2,
  breaklines=true,
  breakatwhitespace=false,
  keywordstyle=\color{blue},
  upquote,
  commentstyle=\color{islamicgreen},
  stringstyle=\color{red},
  morekeywords={*,...},
  frame=none  
         }
        }
      }
{\endtcblisting}

\newenvironment{PC1}[2]{%
\tcblisting{breakable,listing only,colback=cream, enhanced,
 attach boxed title to top left ={yshift=-0.5mm}
    ,title={\href{#1}{ \ttfamily #2}}, sharpish
corners, boxrule=1pt,after, listing options =
   {language=Python,
  basicstyle=\footnotesize\ttfamily,        % the size of the fonts that are used for the code
  numbers=none,                   % where to put the line-numbers
%  numberstyle=\tiny\color{blue},  % the style that is used for the line-numbers
  xleftmargin=-10pt,
  aboveskip=-4pt,
  belowskip=-6pt,
  columns=fixed,
  basewidth=6pt,                     %separation of letters
%  stepnumber=1,                % the step between two line-numbers. If it's 1, each line will be numbered
%  numbersep=7pt,                  % how far the line-numbers are from the code
  backgroundcolor=\color{cream},  % choose the background color. You must add \usepackage{color}
  showspaces=false,               % show spaces adding particular underscores
  showstringspaces=false,         % underline spaces within strings
  showtabs=false,                 % show tabs within strings adding particular underscores
  %frame=single,                   % adds a frame around the code
  %frame=shadowbox,
  %rulecolor=\color{black},        % if not set, the frame-color may be changed on line-breaks within not-black text (e.g. comments (green here))
  tabsize=2,                     % sets default tabsize to 2 spaces
 % captionpos=b,                  % sets the caption-position to bottom
  breaklines=true,                % sets automatic line breaking
  breakatwhitespace=false,        % sets if automatic breaks should only happen at whitespace
  %title=\lstname,                 % show the filename of files included with \lstinputlisting;
                                  % also try caption instead of title
  keywordstyle=\color{blue},      % keyword style
  upquote,     % very important to get right quote for cut/paste
  commentstyle=\color{islamicgreen},   % comment style
  stringstyle=\color{red},      % string literal style
 % escapeinside={\%*}{*)},         % if you want to add a comment within your code
  morekeywords={*,...}
  			 }
    		}
    	}
{\endtcblisting}


\RequirePackage{xcolor}
\definecolor{shadecolor}{rgb}{0.9,0.9,1}

\definecolor{thmcolor}{rgb}{0.97,0.88,0.48}
\definecolor{defncolor}{rgb}{0.76,0.96,0.95}
\definecolor{beaubleu}{rgb}{0.26,0.31,0.61}
\definecolor{darkerblue}{rgb}{0.2,0.2,0.4}
\definecolor{beauvert}{rgb}{0.27,0.52,0.42}


\def\redimensionne{0.76}

\RequirePackage{framed}
\RequirePackage[absolute,overlay]{textpos}
\RequirePackage{graphicx}
\RequirePackage{tikz}
\RequirePackage{calc}

%\newlength{\boxw}
%\newlength{\boxh}
%\newlength{\shadowsize}
%% \newlength{\boxroundness}
%% \newlength{\tmpa}
%% \newsavebox{\shadowblockbox}
%% \setlength{\shadowsize}{6pt}
%% \setlength{\boxroundness}{5pt}
\newcommand{\mathem}[1]{$\boldsymbol{#1}$}

\newcommand{\sca}[1]{\lowercase{#1}}
%\newcommand{\vect}[1]{\lowercase{\boldsymbol{#1}}}
\newcommand{\mat}[1]{\isgreek{#1}\ifthenelse{\boolean{estgrec}}{\matgreektorm{#1}}{\boldsymbol{\mathcal{\uppercase{#1}}}}}



%\tcbuselibrary{listings,skins,breakable,theorems}
\newcommand{\blue}[1]{\textcolor{blue}{#1}}
\newcommand{\m}[1]{\mathbf{#1}}
\renewcommand{\c}[1]{\mathcal{#1}}
%\newcommand{\s}[1]{\mathscr{#1}}
%\newcommand{\idef}{\stackrel{\mathrm{def}}{=}}
%\newcommand{\simiid}{\stackrel{\mathrm{iid}}{\sim}}
%\DeclareMathOperator*{\argmin}{argmin}
%\DeclareMathOperator*{\argmax}{argmax}
\newcommand{\cp}{\stackrel{\bb {P}}{\longrightarrow}}
\newcommand{\cd}{\stackrel{\mathrm{d}}{\longrightarrow}}
\newcommand{\cost}{L}
%\newenvironment{proof}{\paragraph{Proof:}}{\hfill$\square$}

%% \usepackage{algorithm}
%% \usepackage{algorithmic}
%% \algsetup{indent=2em}
%% \graphicspath{{../figures/}}
%% \usepackage{amsmath}
%% \usepackage{amsfonts}
%% \usepackage{amssymb}
%% \usepackage{amsbsy}
%% \usepackage{amsthm}
%% \usepackage{mathrsfs}



\newcommand{\indep}{\rotatebox[origin=c]{90}{$\models$}} % TT: 20/02/18
\newcommand{\loss}{\mathrm{Loss}}
\usepackage{upquote}

\usepackage{tcolorbox}
\tcbuselibrary{listings,skins,breakable,theorems}

\newtcbtheorem[auto counter, number within=section]{thm}{Theorem}
{colback=thmcolor!100,colframe=black!35!black,fonttitle=\bfseries,before
 skip=20pt plus 2pt, before
 skip=20pt plus 2pt}{thm} %"thm" is also the prefix (here) for
                          %referencing, see tcolorbox doc% It has two
                          %arguments: the second is the label

\newtcbtheorem[auto counter, number within=section]{defn}{Definition}
{colback=defncolor,        % Background color for the content
 colframe=headercolor,     % Frame color
 coltitle=white,           % Title color
 fonttitle=\bfseries,      % Title font styling
 colbacktitle=headercolor, % Header background color
 boxed title style={sharp corners=north, rounded corners=south},
 boxrule=0.6mm,            % Slightly thicker frame for better contrast
 before skip=12pt, after skip=12pt, % More consistent spacing
 top=6pt, bottom=6pt, % Controls internal padding
 left=8pt, right=8pt, % Better spacing within the box
 breakable,                % Allows breaking across pages
 enhanced,                 % Enables additional styling
}{defn}

\newtcbtheorem[auto counter, number within=section]{example}{Example}
{colback=examplecolor,        % Background color for the content
 colframe=headercolor,        % Frame color (same as defn)
 coltitle=white,              % Title color
 fonttitle=\bfseries,         % Title font styling
 colbacktitle=headercolor,    % Header background color
 boxed title style={sharp corners=north, rounded corners=south},
 boxrule=0.6mm,               % Slightly thicker frame for better contrast
 before skip=12pt, after skip=12pt, % More consistent spacing
 top=6pt, bottom=6pt, % Controls internal padding
 left=8pt, right=8pt, % Better spacing within the box
 breakable,                 % Allows breaking across pages
 enhanced,                  % Enables additional styling
}{example}

\newtcbtheorem[auto counter, number within=section]{app}{Application}
{colback=appcolor,        % Red content background
 colframe=headercolor,     % Black frame
 coltitle=white,           % White title text
 fonttitle=\bfseries,      % Bold title font
 colbacktitle=headercolor, % Black header background
 boxed title style={sharp corners=north, rounded corners=south},
 boxrule=0.6mm,            % Frame thickness
 before skip=12pt, after skip=12pt, % Spacing
 top=6pt, bottom=6pt,      % Internal padding
 left=8pt, right=8pt,      % Internal spacing
 breakable,                % Page breaks
 enhanced,                 % Extra styling
}{app}

% xattention box
\newtcolorbox{xattention}{breakable, enhanced,
  colback=gray!5, colframe=gray, drop shadow,
  arc=5mm, sharp corners=uphill, % Adjusted corners as per your preference
  boxed title style={boxsep=0pt, right=0pt, left=0pt, toptitle=0pt, colframe=white, colback=white},
  overlay={
    \node[anchor=west, xshift=-12mm, yshift=0cm] at (frame.west)
    {\includegraphics[scale=4.0, keepaspectratio]{warning.png}}; % Use scale instead of width
  },
  boxsep=5pt, before skip=10pt, after skip=10pt
}

% xtip box
\newtcolorbox{xtip}{breakable, enhanced,
  colback=gray!5, colframe=gray, drop shadow,
  arc=5mm, sharp corners=uphill, % Adjusted corners as per your preference
  boxed title style={boxsep=0pt, right=0pt, left=0pt, toptitle=0pt, colframe=white, colback=white},
  overlay={
    \node[anchor=west, xshift=-14mm, yshift=0cm] at (frame.west)
    {\includegraphics[scale=0.25, keepaspectratio]{idea.jpeg}}; % Use scale instead of width
  },
  boxsep=5pt, before skip=10pt, after skip=10pt
}

\newtcolorbox{requiredreading}[1]{ % #1 = title of the box
  breakable,
  enhanced,
  title=#1,                           % Title text
  colback=beauvert!5,                % Green background for the box body
  colframe=headercolor,              % Same frame color as "defn" (dark gray)
  colbacktitle=headercolor,          % Dark header bar
  coltitle=white,                    % White text for the header
  fonttitle=\bfseries,               % Boldface title
  boxed title style={
    sharp corners=north,  % Top corners sharp
    rounded corners=south % Bottom corners rounded
  },
  boxrule=0.5mm,                     % Same thickness as "defn"
  before skip=10pt, 
  after skip=10pt, 
  overlay={
    \node[anchor=west, xshift=-12mm, yshift=0pt] at (frame.west)
    {\includegraphics[scale=0.065]{book.png}}; 
  }
}

\lstset{ 
  language=Matlab,                 % the language of the code
  basicstyle=\footnotesize,        % the size of the fonts that are used for the code
  numbers=left,                    % where to put the line-numbers
  numberstyle=\tiny\color{mygray}, % the style that is used for the line-numbers
  stepnumber=1,                    % the step between two line-numbers. If it's 1, each line will be numbered
  numbersep=5pt,                   % how far the line-numbers are from the code
  backgroundcolor=\color{white},   % choose the background color. You must add \usepackage{color}
  showspaces=false,                % show spaces adding particular underscores
  showstringspaces=false,          % underline spaces within strings
  showtabs=false,                  % show tabs within strings adding particular underscores
  frame=single,                    % adds a frame around the code
  rulecolor=\color{black},         % if not set, the frame-color may be changed on line-breaks within not-black text
  tabsize=2,                       % sets default tabsize to 2 spaces
  captionpos=b,                    % sets the caption-position to bottom
  breaklines=true,                 % sets automatic line breaking
  breakatwhitespace=false,         % sets if automatic breaks should only happen at whitespace
  title=\lstname,                  % show the filename of files included with \lstinputlisting; also try caption instead of title
  keywordstyle=\color{blue},       % keyword style
  commentstyle=\color{mygreen},    % comment style
  stringstyle=\color{mymauve},     % string literal style
  escapeinside={\%*}{*)},          % if you want to add LaTeX within your code
  morekeywords={*,...}             % if you want to add more keywords to the set
}

\addbibresource{references.bib}

\newcommand{\PaperTitle}{CBA Global Markets Research}
\newcommand{\AuthorName}{Zac Kienzle}
\newcommand{\AuthorAffiliation}{Bay Farm Capital}
\newcommand{\AuthorEmail}{zac.kienzle@gmail.com}
\newcommand{\ShortTitle}{Hedging Considerations}

\pagestyle{fancy}
\fancyhf{}
\fancyhead[L]{\nouppercase{\sf \leftmark}}
\fancyhead[R]{\nouppercase{\sf \ShortTitle}}
\fancyfoot[C]{\sf \thepage}
\renewcommand{\headrulewidth}{0.4pt}
\renewcommand{\footrulewidth}{0pt}
\setlength{\headheight}{14.5pt}

\hypersetup{
    hidelinks,
    hypertexnames=false,
    pageanchor=true,
    pdftitle={\PaperTitle},
    pdfauthor={\AuthorName},
    pdfsubject={Quantitative Finance Research},
    pdfkeywords={Market Making, Statistical Arbitrage, Pairs Trading},
    bookmarksopen=true,
    bookmarksnumbered=true,
    pdfpagemode=UseOutlines,
    pdfstartview={FitH},
    unicode=true
}

\title{\PaperTitle}
\author{%
    \AuthorName
    \thanks{\AuthorAffiliation; Email: \href{mailto:\AuthorEmail}{\nolinkurl{\AuthorEmail}}}
}
\date{\today}

\begin{document}

\maketitle
\begin{abstract}
    \noindent
    This is a \LaTeX\ writeup of my progression through a preponderance of alternative hedging strategies
\end{abstract}

\pagenumbering{roman}

\thispagestyle{empty}

\clearpage
\tableofcontents
\listoftables
\listofalgorithms
\thispagestyle{empty}
\clearpage

\pagenumbering{arabic}

\section{Case Brief}

\subsection{Problem Statement}
You and your colleagues work in the Commonwealth Bank of Australia's Global Markets team. Your role is to help clients to navigate their financial market risks. Your client, Prime Property Trust (PPT), is concerned with the potential impact that movements in financial markets may have on their operations. The Treasury team at Prime Property Trust are particularly concerned about interest rate market movements and how it may impact their ability to meet budgeted cost obligations over the next year. The Treasury team wants to know:

\begin{enumerate}
    \item Your view on inflation and the RBA cash rate and; how this will affect Australian interest rates.
    \item The impact of variable interest rates to their business operations and cash flows; and
    \item Financial market instruments that may be appropriate to manage these risks.
\end{enumerate}

\subsection{Background}
Prime Property Trust is an Australian Real Estate Investment Trust (REIT) that manages and invests in a diverse range of property types, including office buildings and retail centres. Prime Property Trust primarily funds its assets with floating rate debt. Prime Property Trust has just announced two new assets to be added to the portfolio:

\begin{enumerate}
    \item Price Corporate Tower (PCT) | | Facility tenor: 3 years | Size of facility: \$100m | Interest rate: 3-month BBSY+1.20\%
    \item Sterling Square (SS) | Facility tenor: 5 years | Size of facility: \$200m | Interest rate: 3-month BBSY+1.50\%
\end{enumerate}

Prime Corporate Tower (PCT) will generate immediate returns and be fully operational from Day 1. Sterling Square (SS) will commence construction Day 1 and be fully operational from FY27 onwards. Assume the PCT facility is refinanced at the end of the 3-year tenor. 

The interest rate on the loan facilities are above and calculated at the beginning of each quarter. Additionally, Prime Property Trust must maintain an Interest Coverage Ratio (ICR) greater than 1.75x at all times or risk breaching the terms of the loan and have its funding withdrawn. Prime Property Trust's cost base and projected revenue for the upcoming financial year is provided on the following page.

\newpage

\subsubsection{Cost Base and Projected Revenue}

\begin{tcolorbox}[
    enhanced,
    breakable,
    title={Prime Property Trust - Income Statement},
    colback=cream,
    colframe=headercolor,
    coltitle=white,
    fonttitle=\bfseries,
    boxed title style={
        sharp corners=north,
        rounded corners=south,
        colframe=headercolor,
        colback=headercolor,
    },
    arc=3mm,
    boxrule=0.5mm,
    before skip=10pt,
    after skip=10pt,
    left=5pt,
    right=5pt,
    top=5pt,
    bottom=5pt,
]
\centering
\textbf{All figures in 000's}
\vspace{5pt}
\footnotesize % Apply smaller font size to the table content
\setlength{\tabcolsep}{2pt} % Reduced column spacing to help fit
\begin{tabular}{l *{10}{r}} % Left aligned for Item, 10 right-aligned for years
\toprule
\textbf{Item} & \textbf{FY24} & \textbf{FY25} & \textbf{FY26} & \textbf{FY27} & \textbf{FY28} & \textbf{FY29} & \textbf{FY30} & \textbf{FY31} & \textbf{FY32} & \textbf{FY33} \\
\midrule
\textbf{Rental Income} &  &  &  &  &  &  &  &  &  &  \\
Retail Rent & 5,850 & 7,116 & 7,289 & 20,057 & 18,542 & 17,126 & 14,496 & 13,046 & 14,095 & 14,971 \\
Office Rent & 3,978 & 2,901 & 2,469 & 5,821 & 7,564 & 9,265 & 12,137 & 14,083 & 14,768 & 13,852 \\
Total Rental Income & 9,828 & 10,017 & 9,758 & 25,878 & 26,106 & 26,391 & 26,633 & 27,129 & 28,863 & 28,823 \\
\midrule
Less Vacancy Factor & - & - & - & - & - & - & - & - & - & - \\
Less Repairs \& Maintenance & 1,047 & 1,078 & 1,176 & 1,232 & 1,309 & 1,424 & 1,518 & 1,607 & 1,760 & 1,925 \\
Less Management Fee & 349 & 527 & 579 & 411 & 436 & 361 & 394 & 536 & 587 & 642 \\
Less Outgoings Paid & - & - & - & - & - & - & - & - & - & - \\
\midrule[\heavyrulewidth]
\textbf{Net Rentals (EBITDA)} & \textbf{8,432} & \textbf{8,412} & \textbf{8,003} & \textbf{24,235} & \textbf{24,361} & \textbf{24,606} & \textbf{24,721} & \textbf{24,986} & \textbf{26,516} & \textbf{26,256} \\
\midrule
\textbf{Base Case} &  &  &  &  &  &  &  &  &  &  \\
BBSY (\%) & 3.70\% & 3.75\% & 3.49\% & 3.36\% & 3.40\% & 3.49\% & 3.59\% & 3.71\% & 3.82\% & 3.92\% \\
\bottomrule
\end{tabular}
\vspace{5pt}
\footnotesize \textit{(Source: Prime Property Trust Income Statement; Assume ICR = EBITDA / Interest Expense)}
\end{tcolorbox}

\newpage
\section{Rates Outlook}

\subsection{Equilibrium and Dynamic Relationships}

\subsection{Macroeconomics}

\newpage

\section{Hedging Objectives}

Maintenance of the debt covenants 

\newpage

\section{Interest Rate Derivatives Valuation}

\newpage

\subsection{Forward Rate Agreements (FRAs)}

A Forward Rate Agreement (FRAs) is an over-the-counter (OTC) contract that locks in an interest rate for a future loan or deposit. The buyer of an FRA locks in the right to borrow at a specific rate, while the seller locks in the right to lend at that rate. The contract is settled with a single cash payment at the settlement date, based on the difference between the agreed-upon forward rate and the actual market interest rate (the "fixing rate") on a notional principal amount. No actual principal is exchanged.

\newpage

\subsection{BAB Futures}

\newpage

\subsection{Interest Rate Swaps}

\newpage

\subsection{Caplets and Floorlets}



\newpage

\section{Hedging Deliberations}

% Quantitative analysis in terms of considerations such as the
% o size or number of contracts to enter into (possibly progressively),
% o position (short or long) and contract maturities,
% o timing and tenor in the market (when and how to enter and exit),
% o initial and ongoing cashflows and costs (premiums, margins, etc),
% o expected outcomes or payoffs under different market conditions,
% o risks and downside if markets didn’t evolve as forecast, of the mechanics each hedging
% option Pershing considered, namely
% not hedging,
% partial or full portfolio liquidation into cash,
% S\&P 500 index futures,
% S\&P 500 index options, and
% index CDS.
% Analysis of the feasibility and suitability of each of Pershing’s hedging options for achieving their
% hedging objectives.

\newpage

\subsection{Not Hedging}
The strategic decision to remain unhedged against interest rate risk implies that PPT consciously accepts the direct financial impact of fluctuations in the 3-month Bank Bill Swap Rate (BBSY). This rate is the floating reference for its newly acquired \$300 million in debt facilities. Such a strategy reflects a managerial assessment where the benefits of potentially favourable interest rate movements outweigh the costs or constraints associated with hedging instruments. This may also indicate that the firm's operational cash flows possess sufficient resilience to absorb interest rate volatility or that hedging instruments' costs are uneconomical (\cite{sr}).
\subsubsection{Mechanics}
The operational mechanics of a non-hedging strategy are characterised by their simplicity, wherein PPT elects to service their debt obligations at the prevailing market rates. Specifically, interest payments on the \$100m Prime Corporate Tower (PCT) facility will be at 3-month BBSY plus a 1.20\% margin, and on the \$200m Sterling Square (SS) facility at 3-month BBSY plus a 1.50\% margin. These interest obligations will reset quarterly, ensuring that the firm's interest expense profile remains directly correlated with short-term interest rate movements.
\subsubsection{Size or Number of Contracts}
Under a strategy of not hedging, PPT does not engage in any off-balance sheet derivative contracts to mitigate interest rate risk. Consequently, the size or number of hedging contracts is zero. The entire nominal value of the floating-rate debt, amounting to \$300 million, remains exposed to interest rate fluctuations.

\subsubsection{Position and Contract Maturities}
Contract maturities are inapplicable in this scenario, as no hedging instruments are utilised. The pertinent maturities are those of the underlying debt facilities: the \$100m PCT facility has a 3-year tenor and is subject to refinancing, while the \$200m SS facility has a 5-year tenor. Therefore, the unhedged interest rate position extends across these periods, with specific managerial concerns directed towards the upcoming 12-month budget cycle.
\subsubsection{Timing and Tenor}
The unhedged strategy is effectuated from the initial drawdown of the debt facilities. It persists for the duration of these loans or until management elects to implement a hedging program. The tenor of this unhedged exposure aligns with the respective tenors of the PCT and SS debt facilities. Each quarterly interest rate reset can be viewed as a discrete point at which the unhedged strategy is implicitly reaffirmed for the subsequent period based on the prevailing BBSY fixing.
\subsubsection{Initial and Ongoing Cashflows and Costs}
No direct initial costs are associated with \emph{not} entering into derivatives contracts to hedge exposures. The primary benefit of not hedging is avoiding upfront premium payments (for options) or potential transaction costs associated with entering other derivatives positions.  

The ongoing cashflows and costs are the variable quarterly interest payments due on the floating rate debt (\$300 million), which is given by:
\begin{equation}
\begin{split}
    \text{Interest Payment} ={}& \left(\text{Principal}_{PCT} \cdot (\text{BBSY}_{3m} + \text{Spread}_{PCT}) \right) /4 \\
    & + \left(\text{Principal}_{SS} \cdot (\text{BBSY}_{3m} + \text{Spread}_{SS}) \right) /4
\end{split}
\end{equation}

The "cost" inherent in this strategy is the economic uncertainty surrounding these payments, which will fluctuate directly with BBSY. Indirect costs may arise while no explicit fees are paid for hedge maintenance. These can include increased managerial resources devoted to monitoring market volatility and re-forecasting and potentially a higher cost of capital or reduced firm valuation if investors perceive the unhedged risk as excessive (\cite{ss}; \cite{sr}).

\subsubsection{Expected Outcomes or Payoffs}
The financial outcome of an unhedged strategy is contingent upon the realised path of future interest rates, juxtaposed with PPT's treasury team's expectations. Contextually, if prevailing market rates rise, PPT weathers escalating interest expenses, which reduce operating cash flows and pressure the ICR, directly eroding firm profitability. Conversely, if interest rates decline, PPT benefits from reduced interest payments, enhancing cash flows and improving the ICR. However, a notable characteristic is the market rate volatility regardless of net direction, wherein PPT's interest expenses become unpredictable, complicating financial planning, budgeting and potentially dividend policy. Such volatility can be particularly detrimental to REITs, which are often valued for stable and predictable income streams. The sensitivity of the annual interest expense and ICR to changes in BBSY is illustrated in Figure \ref{fig:UnhedgedParameterSensitivity}.

\begin{figure}[h] % 'h'ere, 't'op, 'b'ottom, 'p'age of floats
 \centering
 \includegraphics[width=\textwidth]{/Users/zackienzle/CBA/code/plots/Unhedged_Parameter_Sensitivity.png} 
 \caption{Sensitivity of Prime Property Trust's Annual Interest Expense and ICR to BBSY Changes under an Unhedged Strategy.}
    \label{fig:UnhedgedParameterSensitivity}
\end{figure}

The payoff structure is linear concerning BBSY changes; each basis point movement in the 3-month BBSY directly alters the quarterly interest expense by \$7,500 on the combined \$300 million facilities.
\subsubsection{Risks and Downside}
The decision to forego hedging exposes PPT to several material risks amplified by the firm's financial structure and covenants. 

Interest rate risk, the direct risk that increases in the 3-month BBSY, will lead to higher debt servicing costs, given that PPT's explicit communications about budgeted cost obligations constitute a primary risk factor.

Maintaining an interest coverage ratio (ICR) greater than 1.75x is a significant constraint. Previously expounded upon, the pro-forma ICR is approximately $0.54 \text{x}$ (FY25 "Net Rentals (EBITDA)" of \$8,412k and the base case FY25 BBSY of 3.75\%), which is substantially below the covenant threshold. An unhedged strategy offers no protection against further BBSY increases, which would lower the ICR, almost certainly ensuring a breach. The consequences of such a breach can be severe, ranging from the imposition of stricter terms by lenders and demands for early repayment to, ultimately, insolvency (\cite{ss}; \cite{tw}). The high leverage implied by this low ICR makes the firm particularly vulnerable.

Unhedged interest expenses introduce significant volatility into PPT's earnings and cash flow profile, wherein this unpredictability can complicate internal capital budgeting and strategic planning (\cite{fss}), lead to underinvestment if internally generated funds are unexpectedly diverted to higher debt servicing (\cite{ms}) or engender equity market penalisation, especially for REITs where investors often seek stable dividend yields (\cite{ons}).

Firm value considerations.

\subsubsection{Feasibility and Suitability}
The operational feasibility of a non-hedging strategy is indisputable, as it represents the default state requiring no proactive financial market intervention or associated transaction costs. PPT can readily implement this strategy by simply servicing its debt obligations as they fall due, based on the prevailing floating market rates. However, the critical question pertains to the suitability of such an approach for PPT, given its specific financial structure, risk exposures, stated objectives, and the broader expectations for entities within the REIT sector (\cite{mw}).

The decision not to hedge material financial risks can be suboptimal if such risks increase the probability of financial distress or lead to underinvestment in valuable projects (\cite{fss}; \cite{ss}). Moreover, the Treasury team's explicit concern regarding the impact of interest rate movements on their ability to meet budgeted cost obligations over the next year contravenes the inherent uncertainty accepted by forgoing hedging. An unhedged position directly exposes PPT's interest expense and budget adherence to the full volatility of the 3-month BBSY.

The nature of REITs and the expectations of their investors generally favour stable and predictable income streams, often to support regular distributions (\cite{g}). The market can perceive the earnings volatility introduced by unhedged, large-scale floating-rate debt negatively, potentially increasing PPT's cost of equity or reducing its valuation multiples (\cite{aw}). Furthermore, the Sterling Square (SS) development, representing \$200m of the new debt, will be in a non-income-generating construction phase until FY27. During this period, its floating-rate interest payments will be a direct charge against earnings generated by other assets, heightening the sensitivity of PPT's overall financial health to interest rate spikes. Not hedging this specific exposure during the construction phase is particularly risky.

From a managerial decision-making standpoint, the theory of real options might suggest that maintaining flexibility (by not locking into hedges) has value, especially in uncertain environments (\cite{d}). However, this value of flexibility must be weighed against the potentially catastrophic costs of a covenant breach. Given the apparent severity of the ICR situation, the argument for preserving flexibility by not hedging is substantially weakened. The primary objective should shift towards ensuring financial stability and covenant compliance (\cite{t}).

\subsection{Conclusion}
Remaining unhedged, while the simplest option, exposes Prime Property Trust to unmitigated interest rate risk that directly threatens its ability to meet budgeted costs and, critically, to comply with its ICR covenant. Given the projected financials, this strategy is exceptionally high-risk and appears unsuitable for PPT. It would require a radical deviation from prudent financial management unless significant undisclosed information alters the risk landscape.

\newpage

\subsection{Forwards}

\subsubsection{Mechanics}

\subsubsection{Size or Number of Contracts}

\subsubsection{Position and Contract Maturities}

\subsubsection{Timing and Tenor}

\subsubsection{Initial and Ongoing Cashflows and Costs}

\subsubsection{Expected Outcomes or Payoffs}

\subsubsection{Risks and Downside}

\subsubsection{Feasibility and Suitability}

\newpage

\subsection{Futures}
Employing exchange-traded futures contracts represents a standardized approach for PPT to manage its exposure to fluctuating 3-month BBSY rates. 
\subsubsection{Mechanics}
ASX 90-Day Bank Bill Futures are standardized, legally binding agreements to notionally buy or sell a 90-day bank-accepted bill with a face value of AUD $1,000,000$. 
\subsubsection{Size or Number of Contracts}

\begin{equation}
    \text{Number of Contracts} = h^{*} \cdot \frac{\text{Amount to be Hedged}}{\text{Value of One Futures Contract}}
\end{equation}

where $h^{*}$ denotes the hedge ratio. Notably, we can express the hedge ratio as:

\begin{equation}
    h^{*} = \rho_{SF} \frac{\sigma_S}{\sigma_F} = \frac{\text{Cov}(S, F)}{\text{Var}(F)}
\end{equation}
where $S$ is the spot price of the underlying asset, and $F$ is the futures price.

\subsubsection{Position and Contract Maturities}

\subsubsection{Timing and Tenor}

\subsubsection{Initial and Ongoing Cashflows and Costs}

\subsubsection{Expected Outcomes or Payoffs}
The primary expected outcome of employing a short futures hedge is the stabilization of future effective borrowing costs against variability in the 3-month BBSY.


\subsubsection{Risks and Downside}
Substantial adverse movements in futures prices can trigger significant margin calls, requiring PPT to provide liquid funds at short notice. Failure to meet margin calls could lead to the premature and forced liquidation of the hedge at an unfavorable time (\cite{}).

The daily mark-to-market and variation margining process is a significant operational risk.


\subsubsection{Feasibility and Suitability}

\newpage

\subsubsection{Conclusion}

\newpage

\subsection{Swaps}

\subsubsection{Mechanics}
An interest rate swap (IRS) is a contractual agreement between two counterparties to exchange a series of interest payments over a specified period, calculated on a common notional principal amount which itself is typically not exchanged. In its most common form, a "plain vanilla" swap, one party agrees to make payments based on a predetermined fixed interest rate, while the other party agrees to make payments based on a floating interest rate, such as the Bank Bill Swap Rate (BBSY) or another relevant reference rate (REFERENCE, 0000). On each scheduled payment date (e.g., quarterly or semi-annually), these obligations are usually netted, with only the difference being paid by the party owing the larger amount. The fixed rate of the swap is determined at inception such that the initial market value of the swap is typically zero for both parties, reflecting prevailing market conditions and expectations for future interest rates (REFERENCE, 0000). This derivative allows entities to effectively transform the nature of their interest rate exposures or receipts without altering their underlying debt or asset portfolios. 
\subsubsection{Size or Number of Contracts}

\subsubsection{Position and Contract Maturities}
To convert its floating-rate payments into fixed payments, PPT would enter into a payer swap, meaning it agrees to pay the fixed rate and receive the floating rate (3-month BBSY). This effectively neutralizes its floating rate exposure on the debt.

The maturity of the IRS can be customized to match the tenor of the underlying debt exposure PPT wishes to hedge. Specifically, for the 3-year PCT facility, PPT could enter into a 3-year interest rate swap, whereas for the 5-year SS facility, a 5-year swap could be used.

Alternatively, PPT could use shorter-dated swaps and roll them, or use a series of forward-starting swaps to hedge specific future periods, particularly for the concern over the "next year". The ability to precisely match the tenor of the liability is a significant advantage of swaps over futures, which often require rolling for longer-term hedges (REFERENCE, 0000).
\subsubsection{Timing and Tenor}
The swap would ideally be entered into as soon as PPT wishes to lock in a fixed rate for its future interest payments. The fixed rate achievable will depend on the swap market conditions at the time of execution. The "next year" focus () might lead to considering swaps covering that period, or even forward-starting swaps that become effective at a future date.

The tenor of an IRS is highly customizable. PPT can choose tenors ranging from a few months to several years (e.g., 1, 2, 3, 5 years or even longer), allowing for a precise match with the life of its debt facilities or its specific hedging horizon. This alignment reduces the risks associated with maturity mismatches.

\subsubsection{Initial and Ongoing Cashflows and Costs}
Typically, a standard ("at-market") interest rate swap has no upfront premium or initial cost, as the fixed rate is set such that the initial present value of the fixed and expected floating legs are equal. However, transaction costs or bid-ask spreads charged by the swap dealer (the bank) are implicitly built into the agreed fixed rate. Off-market swaps (where the fixed rate is set above or below the prevailing market rate) would involve an upfront payment or receipt.

\subsubsection{Expected Outcomes or Payoffs}
The primary outcome of entering a payer swap is the conversion of PPT's variable interest expense into a predictable, fixed expense for the term of the swap.


\subsubsection{Risks and Downside}

\subsubsection{Feasibility and Suitability}

\newpage

\section{Hedging Rationale}

\newpage

\printbibliography

\end{document}